% GENERATED FILE. Edit canonical lessons, then run npm run generate:books.
% canonical-source-hash: fnv1a64:ef8c46458ebfb3b0
% canonical-lessons: PT-C262-devia, PT-C262-explicar, PT-C262-revisao-explicacao

\chapter{Devia --- What You Should Do, and Explaining Why}
\label{ch:pt-should-and-explaining}
\hlchaptermodality{262}

\begin{chapteropening}
\textbf{By the end of this chapter:} \emph{I can give soft advice with devia and deveria, and explain myself in four moves: an opinion, a cause, a consequence and advice.}

Complete PT-C262-revisao-explicacao: answer a why-question with an opinion, a cause, a consequence and advice.
\end{chapteropening}

\section[devia, deveria]{devia — what you ought to do about it}

\label{lesson:PT-C262-devia}



\begin{quote}
Say \textquotedblleft{}I would like\textquotedblright{}. (\emph{Gostaria.}) Then run a reason forward: \textquotedblleft{}I've got a cold, so I'm not going out.\textquotedblright{} (\emph{Estou constipado, por isso não saio.})

An explanation often ends on a third move: what somebody ought to do. That move has its own form.
\end{quote}

\subsection*{You'll want to know: devia}
You have \textbf{dever}, \textquotedblleft{}to have to, must\textquotedblright{}: \textbf{devo estudar}, I must study. Put it in the imperfeito and it softens into advice:

\begin{quote}
\textbf{Devia} estudar. — I should study. \textbf{Devias} descansar. — You should rest. \emph{(to a friend)} \textbf{Devia} descansar. — You should rest. \emph{(to o senhor, a senhora or você)}
\end{quote}

A bare infinitive follows, with nothing in between. \emph{Devo estudar} has no way out; \emph{devia estudar} says I know I ought to, and have not promised. That softness is the register advice usually wants.

This is the imperfeito doing a present job, and in Portugal it is the everyday way to say \textquotedblleft{}should\textquotedblright{}. The same verb has a second soft form, \textbf{deveria}, with the ending of your \textbf{gostaria}:

\begin{quote}
\textbf{Deveria} descansar mais. — I ought to rest more.
\end{quote}

\emph{Deveria} is a little more careful, and Brazil uses it more than Portugal. Its \emph{-ia} belongs to the would-form of Portuguese, which the last chapter of this run builds properly.

\begin{cousinweb}
\textbf{dever} is inherited from Latin \emph{dēbēre}, \textquotedblleft{}to owe\textquotedblright{} (EtymDB 2.1), and \emph{dēbēre} is a compound: \emph{dē}, \textquotedblleft{}from\textquotedblright{}, plus \emph{habēre}, \textquotedblleft{}to have\textquotedblright{} (Lewis \& Short, under \emph{debeo}). To owe is to \textbf{have something from someone}: you are holding what is theirs. \textbf{a dívida}, \textquotedblleft{}a debt\textquotedblright{}, which you have met, is the same family, and English took it as \textbf{debt}, \textbf{due} and \textbf{duty}.
\end{cousinweb}

\subsection*{Your turn}
Say these aloud:

\begin{itemize}
\raggedright
  \item \emph{devia, devias, devia} — then \emph{deveria}
  \item the hard one, then the soft one — \emph{Devo estudar.} then \emph{Devia estudar.}
  \item advice to a friend with a cold — \emph{Estás constipado, então devias descansar.}
  \item \emph{Acho que devias descansar.}
\end{itemize}

\begin{tcolorbox}[breakable,colback=teal!4,colframe=teal!35!black,title={Before you move on}]
What is the difference between \emph{devo} and \emph{devia}? (\emph{\textbf{Devo}} is \textquotedblleft{}I must\textquotedblright{}; \emph{\textbf{devia}} is \textquotedblleft{}I should\textquotedblright{}, softer.) Which soft form is more usual in Brazil? (\emph{\textbf{Deveria.}}) What are the two Latin pieces of \emph{dēbēre}? (\emph{\textbf{Dē}} and \emph{\textbf{habēre}} — to have from someone.)
\end{tcolorbox}
\section[explicar]{explicar — to unfold}

\label{lesson:PT-C262-explicar}



\begin{quote}
Four pieces are in your hands: \emph{acho que} for an opinion, \emph{porque} and \emph{por isso} for the two arrows, \emph{devia} for advice. This lesson puts them in a row, under the verb that names the act.
\end{quote}

\begin{grammarlens}[title={Grammar Lens: the shape of an explanation}]
\textbf{explicar}, \textquotedblleft{}to explain\textquotedblright{}, is a regular \emph{-ar} verb with one spelling catch. Its event past for \textquotedblleft{}I\textquotedblright{} is \textbf{expliquei}: the \emph{c} becomes \emph{qu} before an \emph{e}, so that it keeps its hard sound. The rest is plain: \emph{explico}, \emph{explicas}, \emph{explica}, \emph{explicou}.

Here is a whole explanation, every word of it taught:

\begin{quote}
— Porque é que não sais? — \textbf{Acho que} vai chover. E estou constipado, \textbf{porque} ontem fazia frio. \textbf{Por isso} não saio. \textbf{Devia} descansar.
\end{quote}

\begin{quote}
\emph{— Why aren't you going out?} \emph{— I think it's going to rain. And I've got a cold, because it was cold yesterday. So I'm not going out. I ought to rest.}
\end{quote}

Four moves, and a word for each: an \textbf{opinion} (\emph{acho que}), a \textbf{cause} pointing back (\emph{porque}), a \textbf{consequence} pointing forward (\emph{por isso}), and \textbf{advice} (\emph{devia}). With both arrows, you can build it cause first or claim first.
\end{grammarlens}

\begin{cousinweb}
\textbf{explicar} is a learned borrowing from Latin \emph{explicāre}, \textquotedblleft{}to unfold, to spread out\textquotedblright{} (EtymDB 2.1; Lewis \& Short, under \emph{explico}), which is \emph{ex-}, \textquotedblleft{}out\textquotedblright{}, plus \emph{plicāre}, \textquotedblleft{}to fold\textquotedblright{}. \textbf{To explain is to unfold}: something was folded shut, and you open it out.

The same \emph{plicāre} also came down into Portuguese the slow way, by mouth, and wore into \textbf{chegar}, \textquotedblleft{}to arrive\textquotedblright{} (EtymDB 2.1): Latin \emph{pl} became Portuguese \emph{ch}. The usual story of how folding became arriving is the sailor's, folding the sails at the end of a voyage; neither source tells it, so it is offered as the usual story only. English kept the Latin in \textbf{explicit}, and \textbf{complicated} is the opposite: folded together.
\end{cousinweb}

\subsection*{Your turn}
Say these aloud:

\begin{itemize}
\raggedright
  \item \emph{explico, explicas, explica} — then \emph{expliquei}
  \item the four moves — \emph{Acho que vai chover, porque não há sol. Por isso não saio. Devia ler.}
  \item the same, consequence first — \emph{Não saio porque chove.}
  \item \emph{explicar}, \emph{chegar}, then English \emph{explicit}
\end{itemize}

\begin{tcolorbox}[breakable,colback=teal!4,colframe=teal!35!black,title={Before you move on}]
Name the four moves of an explanation and a word for each. (Opinion \emph{\textbf{acho que}}, cause \emph{\textbf{porque}}, consequence \emph{\textbf{por isso}}, advice \emph{\textbf{devia}}.) Why is there a \emph{qu} in \emph{expliquei}? (To keep the \emph{c} \textbf{hard}.) What is \emph{explicar}, taken apart? (\emph{\textbf{Ex-}} and \emph{\textbf{plicāre}}: \textbf{to unfold}.)
\end{tcolorbox}
\section[(review)]{Review — explain yourself}

\label{lesson:PT-C262-revisao-explicacao}



\begin{quote}
Say \textquotedblleft{}you should rest\textquotedblright{} to a friend. (\emph{Devias descansar.}) Which English word for owing shares its Latin? (\textbf{Debt}, or \textbf{due}, or \textbf{duty}.)
\end{quote}

\begin{grammarlens}[title={What you've built}]
Two chapters, four moves. You can say what you think, why, what follows from it, and what ought to be done. And because chapter two hundred and fifty-nine gave you the scene, your reasons can reach into the past:

\begin{quote}
— Porque é que não foste ao mercado? — Fazia frio e eu estava constipado, por isso não saí. Acho que devia sair mais.
\end{quote}

The reason is scene, \emph{fazia} and \emph{estava}; the consequence is an event, \emph{não saí}; the advice is about now. Three times in one answer.
\end{grammarlens}

\subsection*{Your turn}
Say these aloud:

\begin{itemize}
\raggedright
  \item the question — \emph{Porque é que não foste ao mercado?}
  \item the answer, all four moves — \emph{Fazia frio e eu estava constipado, por isso não saí. Acho que devia sair mais.}
  \item a different reason — \emph{Chovia, então estudei.}
  \item the same answer, claim first — \emph{Não saí porque fazia frio.}
  \item an explanation of your own: \emph{para mim}, a \emph{porque}, a \emph{por isso}, a \emph{devia}
\end{itemize}

\begin{tcolorbox}[breakable,colback=teal!4,colframe=teal!35!black,title={Before you move on}]
In \emph{fazia frio, por isso não saí}, which half is the scene and which the event? (\emph{\textbf{Fazia frio}} is the scene; \emph{\textbf{não saí}} is the event.) What does \emph{explicar} do, literally? (It \textbf{unfolds}.) And what is softer than \emph{devo}? (\emph{\textbf{Devia.}})
\end{tcolorbox}
